\SourceLeaf{431}%
\[
\begin{aligned}
\frac1r&=A\cos\varphi^{2}+2B\cos\varphi\sin\varphi+C\sin\varphi^{2}\\
&=\frac{A+C}{2}+\frac{A-C}{2}\cos2\varphi+B\sin2\varphi.
\end{aligned}
\]
If we set
\[
\begin{aligned}
\frac{A-C}{2}&=E\cos2\theta\\
B&=E\sin2\theta,
\end{aligned}
\]
where we may take $E$ to have the same sign as $\frac{A-C}{2}$, then
\[
\frac1r=\tfrac12(A+C)+E\cos2(\varphi-\theta).
\]
Clearly $\varphi$ denotes the angle between the cutting plane and another plane passing through the normal and the tangent to which the direction $M$ corresponds. Thus $\frac1r$ plainly has its greatest (absolute) value, or $r$ its least value, when $\varphi=\theta$, and $\frac1r$ has its least value when $\varphi=\theta+90^{0}$. The greatest and least curvatures therefore occur in two planes that intersect at right angles. The extreme values of $\frac1r$ are consequently
\[
\tfrac12(A+C)\pm\BareRoot{\left\{\left(\frac{A-C}{2}\right)^{2}+BB\right\}},
\]
their sum is $=A+C$ and their product $=AC-BB$, or the product of the two extreme radii of curvature is
\[
=\frac{1}{AC-BB}.
\]
This product is of great importance and deserves to be worked out in greater detail. From the formulas above we find
\[
AC-BB=(\alpha\beta'-\beta\alpha')^{2}(TV-UU)ZZ.
\]
From the third formula in [theorem] $6$, Art. $7$, however, one readily infers that
\[
\alpha\beta'-\beta\alpha'=\pm Z,
\]
hence
\[
AC-BB=Z^{4}(TV-UU).
\]
Furthermore, by Art. $8$ we have
